\section{GİRİŞ}
\label{bolum1}

\subsection{Problem Tanımı}
\label{problem-tan-m}

Modern web tabanlı uygulamaların büyük çoğunluğu, istemci ve sunucu arasındaki veri iletişimini REST API mimarisi üzerinden gerçekleştirmektedir. Günümüzde mobil uygulamalar, mikro servis sistemleri, IoT platformları, bulut tabanlı servisler ve modern web uygulamaları; veri alışverişi, kullanıcı doğrulama, servis entegrasyonu ve sistemler arası haberleşme işlemleri için yoğun biçimde API teknolojilerine ihtiyaç duymaktadır. REST mimarisi; hafif yapısı, platform bağımsız çalışabilmesi, HTTP protokolü ile uyumlu olması ve ölçeklenebilirliği nedeniyle günümüzde en yaygın kullanılan servis mimarilerinden biri haline gelmiştir. Özellikle frontend ve backend katmanlarının birbirinden ayrıldığı modern yazılım mimarilerinde REST API’ler, sistemin temel iletişim omurgasını oluşturmaktadır.

REST API kullanımının yaygınlaşmasıyla birlikte güvenlik problemleri de önemli ölçüde artmıştır. API’ler doğrudan internet üzerinden erişilebilir olduğu için saldırganlar açısından önemli hedefler haline gelmiştir. Yetkisiz erişim, SQL Injection saldırıları, Cross-Site Scripting (XSS), brute force giriş denemeleri, token manipülasyonu, Broken Authentication, güvenlik doğrulama eksiklikleri, hatalı yetkilendirme mekanizmaları ve servis yoğunlaştırma (DoS/DDoS) saldırıları gibi tehditler; güvenli olmayan API sistemlerinde ciddi veri kayıplarına, kullanıcı bilgilerinin sızdırılmasına, servis kesintilerine ve sistem bütünlüğünün bozulmasına neden olabilmektedir. Özellikle kullanıcı doğrulama ve veri erişim kontrollerinin yanlış yapılandırıldığı API sistemlerinde saldırganlar, sistem içerisindeki hassas verilere erişebilmekte veya sistem davranışını manipüle edebilmektedir.

OWASP (Open Web Application Security Project) tarafından yayınlanan “API Security Top 10” listesi, modern API sistemlerinde en sık karşılaşılan güvenlik açıklarını ve bu açıkların ne kadar kritik sonuçlar doğurabileceğini açıkça göstermektedir. Bu listede yer alan Broken Object Level Authorization (BOLA), Broken Authentication, Excessive Data Exposure, Security Misconfiguration ve Injection saldırıları; günümüzde birçok gerçek sistemde karşılaşılan temel güvenlik problemleri arasında yer almaktadır. Özellikle hızlı geliştirme süreçlerinde güvenlik kontrollerinin ikinci plana atılması, geliştiricilerin yalnızca işlevselliğe odaklanması ve güvenlik testlerinin yeterince yapılmaması; REST API tabanlı sistemlerde ciddi zafiyetlerin oluşmasına neden olmaktadır.

Bununla birlikte REST API güvenliği konusunda mevcut eğitim materyallerinin büyük çoğunluğu teorik düzeyde kalmakta; geliştiricilerin gerçek saldırı ve savunma mekanizmalarını deneyimleyebileceği kontrollü laboratuvar ortamları yeterince yaygın bulunmamaktadır. Birçok eğitim içeriği yalnızca kavramsal açıklamalar sunmakta, ancak güvenlik açıklarının uygulama üzerinde nasıl oluştuğunu, saldırıların sisteme nasıl etki ettiğini ve bu saldırılara karşı hangi koruma mekanizmalarının nasıl çalıştığını pratik olarak göstermemektedir. Bu durum, güvenli yazılım geliştirme farkındalığının artırılmasında önemli bir eksikliğe yol açmaktadır.

Ayrıca geliştiricilerin önemli bir kısmı, REST API güvenliğini yalnızca kullanıcı giriş sistemi veya token doğrulama mekanizması olarak değerlendirmekte; input validation, rate limiting, güvenli hata yönetimi, erişim kontrolü, veri filtreleme ve saldırı tespit mekanizmaları gibi kritik güvenlik katmanlarını yeterince uygulamamaktadır. Özellikle küçük ve orta ölçekli projelerde güvenlik önlemleri çoğu zaman proje geliştirme sürecinin son aşamasına bırakılmakta veya tamamen göz ardı edilmektedir. Bu durum, saldırganların sistemlere çok daha kolay erişebilmesine zemin hazırlamaktadır.

Bu proje kapsamında geliştirilen güvenli REST API laboratuvar sistemi ile; hem modern REST API mimarisinin temel çalışma prensiplerinin gösterilmesi hem de yaygın güvenlik açıklarının kontrollü bir ortamda analiz edilmesi amaçlanmıştır. Sistem içerisinde kullanıcı doğrulama mekanizmaları, modüler API yapısı, rate limiting, regex tabanlı temel WAF kontrolleri, CORS güvenliği ve saldırı filtreleme mekanizmaları uygulanarak güvenli backend geliştirme yaklaşımı örneklenmiştir. Böylece geliştiricilerin yalnızca teorik bilgi edinmesi değil, aynı zamanda gerçekçi saldırı senaryolarını ve savunma mekanizmalarını uygulamalı olarak deneyimleyebilmesi hedeflenmiştir.

Bu çalışma aynı zamanda güvenli yazılım geliştirme kültürünün yaygınlaştırılmasına katkı sağlamayı amaçlamaktadır. Projede geliştirilen yapı sayesinde geliştiriciler; REST API sistemlerinde oluşabilecek güvenlik açıklarını analiz edebilmekte, saldırı yüzeylerini gözlemleyebilmekte ve modern backend mimarilerinde uygulanması gereken temel güvenlik prensiplerini pratik olarak inceleyebilmektedir. Böylece proje, hem eğitimsel hem de teknik açıdan REST API güvenliği konusunda uygulamalı bir öğrenme ortamı sunmaktadır.

\subsection{Proje Amaçları}
\label{proje-ama-lar}

Bu çalışmanın temel amacı, REST API güvenlik açıklarının teorik bilgilerden öte uygulamalı olarak öğrenilmesini sağlayan eğitim tabanlı bir saldırı ve savunma simülasyon laboratuvarı geliştirmektir. Günümüzde REST API teknolojileri modern yazılım sistemlerinin temel bileşenlerinden biri haline gelmiş olsa da API güvenliği konusunda uygulamalı eğitim ortamlarının yetersiz olması, geliştiricilerin gerçek saldırı senaryolarını yeterince deneyimleyememesine neden olmaktadır. Bu proje ile birlikte hem saldırı tekniklerinin hem de modern savunma mekanizmalarının kontrollü bir laboratuvar ortamında deneyimlenebilmesi hedeflenmiştir.

Çalışmanın bir diğer amacı, güvenli yazılım geliştirme süreçlerinde karşılaşılan yaygın güvenlik problemlerinin gerçek sistem mimarileri üzerinde analiz edilmesini sağlamaktır. Bu kapsamda geliştirilen sistem yalnızca teorik bilgi sunan bir eğitim platformu değil; aynı zamanda kullanıcıların güvenlik açıklarını keşfedebildiği, saldırı mantığını anlayabildiği ve savunma mekanizmalarının sistem davranışını nasıl değiştirdiğini gözlemleyebildiği uygulamalı bir siber güvenlik laboratuvarı olarak tasarlanmıştır.

Bu proje kapsamında OWASP tarafından yayınlanan API Security Top 10 listesindeki kritik güvenlik açıklarının gerçek bir sistem üzerinde simüle edilmesi amaçlanmıştır. Özellikle Broken Authentication, Broken Object Level Authorization (BOLA), Injection saldırıları, Security Misconfiguration, Excessive Data Exposure ve Rate Limiting eksiklikleri gibi yaygın API zafiyetlerinin hem saldırı hem de savunma perspektifinden incelenmesi hedeflenmiştir. Böylece kullanıcıların yalnızca güvenlik açığının teorik tanımını öğrenmesi değil; bu açığın sistem içerisinde nasıl oluştuğunu, saldırgan tarafından nasıl kullanılabildiğini ve hangi savunma yöntemleriyle engellenebileceğini uygulamalı olarak deneyimlemesi amaçlanmıştır.

Projenin önemli hedeflerinden biri de öğrencilere ve geliştiricilere CTF (Capture The Flag) mantığına benzer etkileşimli bir öğrenme deneyimi sunmaktır. Bu yapı sayesinde kullanıcılar yalnızca pasif şekilde eğitim materyali okumak yerine; gerçek saldırı senaryolarını çözmeye çalışarak sistem güvenliği konusunda aktif deneyim kazanabilmektedir. Özellikle güvenlik açıklarının kontrollü görevler halinde sunulması, kullanıcıların saldırı mantığını adım adım analiz edebilmesine olanak sağlamaktadır. Böylece öğrenme süreci daha interaktif, kalıcı ve uygulama odaklı hale getirilmektedir.

Çalışmanın bir diğer temel amacı, aynı sistem üzerinde hem güvensiz (vulnerable) hem de güvenli (secure) API ortamları oluşturarak karşılaştırmalı analiz yapılabilmesini sağlamaktır. Güvensiz ortamda kullanıcılar çeşitli güvenlik açıklarını deneyimleyebilirken; güvenli ortamda aynı saldırıların neden başarısız olduğu gözlemlenebilmektedir. Bu yaklaşım sayesinde modern savunma mekanizmalarının sistem güvenliği üzerindeki etkisi doğrudan analiz edilebilmektedir. Özellikle güvenlik katmanlarının etkinliğini uygulamalı olarak görmek, geliştiricilerin güvenli backend mimarisi konusundaki farkındalığını artırmaktadır.

Proje kapsamında modern API güvenlik mekanizmalarının uygulamalı olarak gösterilmesi de hedeflenmiştir. Bu doğrultuda JWT (JSON Web Token) tabanlı kimlik doğrulama sistemi, rol tabanlı erişim kontrolü (RBAC), regex tabanlı temel WAF (Web Application Firewall) filtreleme mekanizması, rate limiting sistemi ve güvenli yapılandırma kontrolleri sisteme entegre edilmiştir. Böylece kullanıcılar yalnızca saldırı tekniklerini değil; aynı zamanda modern backend sistemlerinde kullanılan profesyonel savunma yöntemlerini de inceleyebilmektedir. Özellikle JWT yapısının çalışma mantığı, token doğrulama süreçleri, erişim kontrol mekanizmaları ve istek filtreleme sistemleri proje içerisinde uygulamalı olarak gösterilmektedir.

Bu çalışmanın eğitimsel hedeflerinden biri de güvenli yazılım geliştirme kültürünün yaygınlaştırılmasına katkı sağlamaktır. Günümüzde birçok yazılım projesinde güvenlik önlemleri geliştirme sürecinin son aşamasında ele alınmakta veya tamamen göz ardı edilmektedir. Bu proje ile birlikte güvenliğin yazılım geliştirme sürecinin merkezinde yer alması gerektiği vurgulanmaktadır. Özellikle geliştiricilerin daha proje tasarım aşamasında güvenlik odaklı düşünme alışkanlığı kazanması hedeflenmiştir.

Ayrıca proje, siber güvenlik eğitimi alan öğrenciler için uygulamalı bir laboratuvar ortamı sunmayı amaçlamaktadır. Kullanıcılar bu sistem sayesinde API güvenliği, kimlik doğrulama mekanizmaları, erişim kontrol sistemleri, saldırı tespit yöntemleri ve güvenlik filtreleme teknikleri gibi konuları gerçekçi bir backend mimarisi üzerinde deneyimleyebilmektedir. Böylece proje yalnızca teknik bir yazılım geliştirme çalışması değil; aynı zamanda eğitim ve farkındalık odaklı bir siber güvenlik platformu niteliği taşımaktadır.

Sonuç olarak bu çalışmanın temel amacı; REST API güvenliği konusunda teorik bilgi ile pratik uygulama arasındaki boşluğu azaltmak, geliştiricilerin gerçek dünya saldırı senaryolarını kontrollü şekilde deneyimleyebilmesini sağlamak ve modern güvenlik mekanizmalarının çalışma mantığını uygulamalı olarak öğretmektir. Bu sayede proje, hem yazılım geliştirme hem de siber güvenlik eğitimi açısından kapsamlı bir öğrenme ve analiz ortamı sunmaktadır.

\subsection{Proje Kapsamı}
\label{proje-kapsam}

Bu proje, modern REST API güvenlik yaklaşımlarının uygulamalı olarak analiz edilebilmesi amacıyla geliştirilmiş, Docker tabanlı modüler bir siber güvenlik laboratuvar sistemidir. Sistem, aynı konteyner ortamı içerisinde çalışan ve birbirinden bağımsız şekilde yapılandırılmış iki ayrı REST API backend yapısından oluşmaktadır. Bu backend’lerden ilki kasıtlı olarak güvenlik açıkları içeren “güvensiz API ortamını (vulnerable environment)”, ikinci backend ise modern savunma mekanizmalarıyla güçlendirilmiş “güvenli API ortamını (secure environment)” temsil etmektedir. Böylece kullanıcılar aynı işlevleri yerine getiren iki farklı API mimarisini karşılaştırmalı olarak inceleyebilmekte ve güvenlik mekanizmalarının sistem davranışı üzerindeki etkisini uygulamalı biçimde gözlemleyebilmektedir.

Proje kapsamında geliştirilen her iki backend sistemi de gerçek dünya REST API mimarisine benzer şekilde tasarlanmıştır. Sistemlerde kullanıcı yönetimi, kullanıcı kayıt ve giriş işlemleri, JWT tabanlı kimlik doğrulama, rol tabanlı erişim kontrolü (RBAC), veri listeleme, veri oluşturma, güncelleme ve silme gibi temel API işlemleri bulunmaktadır. Böylece proje yalnızca teorik saldırı örnekleri sunan basit bir demo sistemi olmaktan çıkarılarak gerçekçi backend mimarisine yakın bir laboratuvar ortamı haline getirilmiştir.

Güvensiz API ortamında çeşitli güvenlik açıkları bilinçli şekilde sisteme entegre edilmiştir. Bu ortamda kullanıcılar SQL Injection, zayıf kimlik doğrulama mekanizmaları, eksik erişim kontrolü, yetersiz input validation, Broken Object Level Authorization (BOLA), hatalı güvenlik yapılandırmaları ve rate limiting eksiklikleri gibi yaygın REST API zafiyetlerini analiz edebilmektedir. Bu yapı sayesinde saldırgan bakış açısıyla sistem davranışı incelenebilmekte ve güvenlik açıklarının oluşma nedenleri uygulamalı olarak gözlemlenebilmektedir.

Güvenli API ortamında ise aynı işlevler modern güvenlik prensiplerine uygun şekilde yeniden yapılandırılmıştır. Bu ortamda JWT tabanlı güvenli kimlik doğrulama sistemi, rol tabanlı erişim kontrolü, regex tabanlı temel WAF (Web Application Firewall) filtreleme sistemi, rate limiting mekanizması, güvenli input validation işlemleri, güvenli hata yönetimi ve CORS güvenlik yapılandırmaları uygulanmıştır. Böylece kullanıcılar yalnızca saldırıların nasıl gerçekleştiğini değil; aynı zamanda bu saldırıların hangi savunma yöntemleriyle engellenebileceğini de doğrudan analiz edebilmektedir.

Proje kapsamında Docker teknolojisi kullanılarak tüm sistemin konteyner tabanlı çalışması sağlanmıştır. Bu yaklaşım sayesinde hem güvenli hem de güvensiz backend servisleri izole şekilde çalıştırılabilmekte, sistemin kurulumu kolaylaşmakta ve laboratuvar ortamının taşınabilirliği artırılmaktadır. Docker kullanımı aynı zamanda gerçek dünya mikro servis mimarilerine benzer bir geliştirme ortamı sunmaktadır. Böylece kullanıcılar yalnızca API güvenliği değil, modern backend dağıtım mantığı hakkında da deneyim kazanabilmektedir.

Sistem içerisinde frontend tarafı temel düzeyde tutulmuş, ana odak backend güvenliği ve API davranışlarının analiz edilmesine verilmiştir. Bu nedenle proje kapsamı içerisinde gelişmiş kullanıcı arayüzü geliştirme süreçleri yer almamaktadır. Kullanıcı etkileşimleri ağırlıklı olarak REST istemcileri, HTTP istekleri ve API endpoint testleri üzerinden gerçekleştirilmektedir. Böylece dikkat doğrudan güvenlik mekanizmaları ve backend davranışlarına yönlendirilmiştir.

Proje kapsamında özellikle uygulama katmanı güvenliği hedeflenmiş, ağ seviyesindeki gelişmiş saldırı senaryoları kapsam dışında bırakılmıştır. Distributed Denial of Service (DDoS), Man-in-the-Middle (MITM), DNS spoofing, packet sniffing ve düşük seviyeli ağ saldırıları gibi konular bu çalışmanın sınırları dışında tutulmuştur. Bunun temel nedeni, çalışmanın odak noktasının ağ güvenliği değil REST API uygulama güvenliği olmasıdır. Ayrıca bu tür saldırıların analiz edilebilmesi için daha karmaşık ağ topolojileri, sanal ağ yapılandırmaları ve gelişmiş güvenlik altyapıları gerekmektedir.

Bununla birlikte proje kapsamında makine öğrenimi tabanlı saldırı tespit sistemleri veya anomali analiz mekanizmaları da yer almamaktadır. Günümüzde birçok modern güvenlik sistemi yapay zekâ destekli saldırı tespiti, trafik analizi ve davranış modelleme yöntemleri kullanmasına rağmen bu çalışma temel olarak REST API güvenlik mimarisi ve uygulama katmanı savunmaları üzerine odaklanmıştır. Ancak ilerleyen çalışmalarda sisteme yapay zekâ tabanlı anomali tespiti, saldırı davranışı analizi ve otomatik tehdit sınıflandırma mekanizmalarının eklenmesi planlanmaktadır.

Ayrıca proje kapsamı içerisinde gerçek üretim ortamlarında kullanılan gelişmiş güvenlik servisleri, merkezi log yönetimi, SIEM entegrasyonları, dağıtık güvenlik izleme sistemleri ve kurumsal seviyede API gateway çözümleri yer almamaktadır. Bu sistemler gelecekte projenin daha ileri sürümlerinde genişletilebilir modüller olarak değerlendirilmektedir.

Sonuç olarak bu proje; REST API güvenliğinin uygulamalı olarak öğrenilebilmesi amacıyla geliştirilmiş, güvenli ve güvensiz backend sistemlerini karşılaştırmalı şekilde sunan Docker tabanlı bir eğitim ve analiz laboratuvarı niteliğindedir. Proje kapsamında temel odak uygulama katmanı güvenliği, saldırı-senaryo analizi ve modern savunma mekanizmalarının incelenmesi olurken; gelişmiş ağ saldırıları, frontend odaklı kullanıcı deneyimi geliştirmeleri ve yapay zekâ destekli güvenlik sistemleri kapsam dışında bırakılmıştır. Böylece proje, kontrollü ve anlaşılabilir bir laboratuvar ortamında REST API güvenlik farkındalığını artırmayı amaçlayan kapsamlı bir uygulama platformu sunmaktadır.


\clearpage
